% !TITLE 诗经·秦风·蒹葭
% !AUTHOR 无名氏
% !DYNASTY 先秦
% !GENRE 四言诗
% !ORDER 7

\begin{poem}
蒹葭\textsuperscript{1}苍苍，白露为霜。
所谓伊人\textsuperscript{2}，在水一方。
溯洄\textsuperscript{3}从之，道阻且长。
溯游从之，宛在水中央。

蒹葭萋萋，白露未晞\textsuperscript{4}。
所谓伊人，在水之湄\textsuperscript{5}。
溯洄从之，道阻且跻\textsuperscript{6}。
溯游从之，宛在水中坻\textsuperscript{7}。

蒹葭采采，白露未已。
所谓伊人，在水之涘。
溯洄从之，道阻且右。
溯游从之，宛在水中沚。
\end{poem}

\begin{annotations}
\notetext{1}{蒹葭：芦苇。蒹，没长花穗的芦苇；葭，初生的芦苇。苍苍、萋萋、采采都形容其茂盛。}
\notetext{2}{伊人：那个人，指心中爱慕或追寻的对象。}
\notetext{3}{溯洄、溯游：溯洄，逆流而上；溯游，顺流而下。从，追寻。}
\notetext{4}{晞：晞（xī），干、晒干，指露水还没干。}
\notetext{5}{湄、涘：湄（méi），水草与水相接之处；涘（sì），水边、岸边。}
\notetext{6}{跻：跻（jī），升高，引申为道路陡峭难行。右，弯曲，指道路曲折。}
\notetext{7}{坻、沚：坻（chí），水中的沙洲、小高地；沚（zhǐ），水中极小的沙洲。}
\end{annotations}

\begin{appreciation}
《蒹葭》是《诗经》中最富有意境美的诗篇之一，也是中国文学史上最早、最动人的"追寻"主题。

"蒹葭"（jiān jiā），就是芦苇。苍苍、萋萋、采采，都是形容芦苇茂盛的样子，但程度略有不同：苍苍是深青色，带有秋天的肃杀；萋萋是浅绿中带着生机；采采是颜色鲜明、长势旺盛。三章分别对应秋天的三个阶段：白露为霜（已经结霜）、白露未晞（露水还没干）、白露未已（露水还没消失）。时间在推移，秋天的寒意越来越深，但追寻者却始终没有放弃。

"所谓伊人，在水一方"——我心中的那个人，在水的另一边。"伊人"是谁？可以是恋人，可以是贤人，可以是理想，也可以是一种可望而不可即的美好事物。这首诗的妙处，就在于它的模糊性——它不告诉你追寻的对象是什么，只告诉你追寻本身的状态。

"溯洄从之，道阻且长"——逆流而上追寻她，道路艰险又漫长。"溯游从之，宛在水中央"——顺流而下追寻她，她仿佛就在水中央。但无论怎么走，都到不了。她好像在那里，又好像不在；好像很近，又好像很远。这种若即若离的感觉，是所有追寻者的共同体验。

三章的结构几乎相同，只换了几个字，却造成了回环往复的音乐效果。这种手法叫"重章叠句"，是《诗经》最典型的艺术特征。它不像现代诗那样追求变化，而是在重复中深化情感。每重复一次，追寻的执着就加深一层，秋天的寒意就浓重一分，那种可望不可即的怅惘就更强烈一些。

这首诗对后世的影响极大。"在水一方"成为中国文学中"可望而不可即"的经典意象。王国维在《人间词话》里说它"最得风人深致"，意思是它最能体现《诗经》那种含蓄深远的韵味。读这首诗，你会感到一种淡淡的忧伤——不是撕心裂肺的痛苦，而是一种更持久、更普遍的惆怅：人生中总有一些美好的东西，我们永远追求，却永远得不到。
\end{appreciation}
